\enteteExam{NFP 107 -- Systèmes de gestion de bases de données - Durée 2h30}%
{Examen première session 11 juin 2025}%
{Examen sur table -- Seuls documents autorisés\,: 10 pages maximum de notes.}%


On met en place un système d'information sur les structures administratives françaises\,: 
régions, départements, préfectures.
Voici le schéma de la base sur laquelle nous allons travailler. 
Les clés primaires sont en gras, les clés étrangères ne sont pas indiquées.

\begin{itemize}
\item Personne(\textbf{idPersonne}, nom, prénom)
\item Région(\textbf{codeRégion},  intitulé, préfecture, idPrésident)
\item Département (\textbf{codeDpt}, nom, préfecture, population, codeRégion, idPrésident)
\item Voisins (\textbf{codeDpt1}, \textbf{codeDpt2})
\end{itemize}

Dans les tables \texttt{Région} et  \texttt{Département}, \texttt{préfecture} désigne
la ville siège de la préfecture de région ou de département. Strasbourg
est à la fois préfecture de la région \textit{Grand Est} et du  département
du Bas-Rhin. Dans la même région, Colmar est préfecture du département du Haut-Rhin.
Régions et départements sont présidés par une personne.

Deux départements sont voisins s'ils ont une frontière commune.

\section{Compréhension du schéma (6 pts)}

\begin{enumerate}
 \item Donner les commandes de création de table pour \texttt{Département} et \texttt{Voisins}.
       On supposera que \texttt{Personne} et \texttt{Région} ont déjà été créées.
  \item On ajoute la commande suivante au schéma de la table \texttt{Voisins}
  
\begin{minted}{sql}
check (codeDpt1  < codeDpt2)
\end{minted}
  
   Quelle est à votre avis la motivation de cette contrainte\,?

  \item Donner le schéma entité-association 
        correspondant à ce schéma relationnel.
        
   \item Existe-t-il une association plusieurs-plusieurs dans ce schéma\,? Si oui 
     vous paraîtrait-il judicieux de la réifier\,? Justifiez votre réponse.

  \item Une même personne peut-elle présider plusieurs départements\,? 
   Peut-elle présider un département et une région\,?
   \item Notons \texttt{D} le département, \texttt{P} la personne, 
   \texttt{ND} le nom du département et \texttt{NP} le nom de la personne.
   Le schéma nout donne les dépendances fonctionnelles
   suivantes\,: \texttt{D -> P, ND} et \texttt{P -> NP}.
   
   On ajoute la dépendance \texttt{P -> D}. Comment l'interpréter
   et que devient le schéma relationnel\,?
   
\end{enumerate}

Dans toute la suite de l'examen, on travaille sur 
le schéma donné dans l'énoncé, et en supposant que la contrainte
sur l'ordre des codes de département est respectée dans la table \texttt{Voisins}
(donc,
\texttt{codeDpt1  $<$ codeDpt2}).

\begin{solution}
 \begin{enumerate}
 
      \item Standard. Voici pour la table \texttt{Voisins}:
     
\begin{minted}{sql}
create table Voisins (idDpt1 int not null, 
    idDpt2 int not null
     primary key (idDpt1, idDpt2),
     foreign key (idDpt1) references Département (idDpt),
         foreign key (idDpt2) references Département (idDpt)
     )
\end{minted}
  
    \item Il existe une symétrie dans la notion de voisin\,: si A
    est voisin avec B, B est voisin avec A. Pour éviter le doublon
   de représentation (voisin(A, B) et voisin(B,A)), on décide 
   de n'accepter que le premier, d'où la clause \texttt{check} 
     \item 
     
     \item \texttt{Voisins} provient d'une association plusieurs-plusieurs.
     La contrainte sur la clé implique qu'un département ne peut pas
     être deux fois le voisin d'un autre. C'est ici tout à fait raisonnable,
     donc la réification est inutile.
 \end{enumerate}
 
\end{solution}



\section{Algèbre et vues (4 pts)}

On définit algébriquement les relations suivantes:

\begin{itemize}
  \item $V1 := Voisins$
  \item $V2 := \rho_{codeDpt1 \to codeDpt2, codeDpt2 \to codeDpt1} (Voisins)$
  \item $V3 := V1 \cup V2$
\end{itemize}

\begin{enumerate}
  \item Quel est le résultat de $V1 \cap V2$\,?
  \item Comment caractériseriez-vous le contenu  de $V3$\,?
    \item Donnez la commande SQL de création de la vue $V3$.
    \item Soit le département $d_1$  de code 
      $c_1$. Quelle requête algébrique sur $V3$ donne les codes de \emph{tous} ses voisins
      (expliquer). 
\end{enumerate}


\begin{solution}

\begin{enumerate} 
  \item Le résultat est vide\,!

  \item Dans $V2$ on inverse la relation entre les codes de département. Donc, dans V2,
     \texttt{codeDpt2  $<$ codeDpt1}). Du coup, dans V3, chaque relation de voisinage
     entre deux départements est représentée dans les deux sens.
  
  \item On peut donc créer la vue suivante qui pourra nous simplifier les requêtes SQL
     par la suite.
     
\begin{minted}{sql}
create view V3 as
select *
from Voisins 
union
select codeDpt2 as codeDpt1, codeDpt1 as codeDpt2
from Voisins
\end{minted}

   \item Avec V3, on n'a plus besoin de se soucier de l'ordre des codes de département.
     On peut donc écrire soit\,:
   
   $$\pi_{codeDpt2} (\sigma_{codeDpt1=c_1} (\rm{V3}))$$   
   ou
   $$\pi_{codeDpt1} (\sigma_{codeDpt2=c_1} (\rm{V3}))$$
   qui donneront le même résultat.

\end{enumerate}

\end{solution}


\section{SQL (7 pts)}

Exprimez en SQL les requêtes suivantes\,:

\begin{enumerate}
  \item Qui préside (prénom, nom) la région contenant le département 'Cantal'\,?
  
  \item Quelles villes sont à la fois préfecture de région et d'un 
  département de plus de 100\,000 habitants
  (donner le nom de la ville, du département, et l'intitulé de la région). 
  
  \item Quels départements d'une même région sont voisins l'un de l'autre\,?
     Donnez l'intitulé de la région et les noms des deux départements.
     Attention à bien prendre en compte la contrainte sur les clés
     dans \texttt{Voisins}
     
% \item Donner les paires de régions ayant une frontière commune
%   (nom des régions). 
%   
%   Aide\,: et si vous utilisiez la vue V3 définie précédemment\,?
   
   \item Quelles régions n'ont pas de département\,?
   
   \item Quelles personnes ne président ni région ni département\,?
   
   \item Donner, pour chaque région, le nombre de département et sa population totale (obtenue
   par cumul de celle des départements)
   
   \item Dans la requête précédente, que se passe-t-il si une
      région n'a pas de département\,? Comment réussir 
      à afficher le nom de la région avec la valeur 0 pour le nombre
      de départements dans ce cas\,?
\end{enumerate}


\begin{solution}


\begin{enumerate}
  \item 
\begin{minted}{sql}
select prénom, nom
from Personne as p, Région as r, Département as d
where d.nom='Cantal'
and d.codeRégion = r.codeRégion
and r.idPrésident = p.idPersonne
\end{minted}

  \item 
\begin{minted}{sql}
select r.préfecture as ville, r.intitulé, d.noms
from  Région as r, Département as d
where r.préfecture = d.préfecture
and d.population > 10000
\end{minted}


  \item 
\begin{minted}{sql}
select r.intitulé, d1.nom as nomDpt1, d2.nom as nomDpt2, Voisins as v
from  Région as r, Département as d1, Département as d2
where r.codeRégion = d1.codeRégion
and r.codeRégion = d2.codeRégion
/* Pour prendre la relation 'voisins dans le bon sens' */
and d1.codeDpt < d2.codeDpt
and v.codeDpt1=d1.codeDpt
and v.codeDpt2=d2.codeDpt
\end{minted}

\begin{comment}
\item Solution avec la vue V3.
\begin{minted}{sql}
select r.intitulé, d1.nom as nomDpt1, d2.nom as nomDpt2
from  Région as r1, Région as r2, Département as d1, Département as d2, V3 as v
where d1.codeRégion = r1.codeRégion
and d2.codeRégion = r2.codeRégion
/* On sait que dans V3, on trouve chaque voisinage représenté dans les deux sens */
and v.codeDpt1=d1.codeDpt
and v.codeDpt2=d2.codeDpt
and r1.codeRégion != r2.codeRégion
\end{minted}

Sinon, solution SQL complète avec un 'ou' logique
\begin{minted}{sql}
select r1.intitulé as nomRégion1, r2.intitulé as nomRégion2
from  Région as r1, Région as r2, Département as d1, 
          Département as d2, Voisins as v
where d1.codeRégion = r1.codeRégion
and d2.codeRégion = r2.codeRégion
and ( 
      (d1.codeDpt < d2.codeDpt and v.codeDpt1 = d1.codeDpt1 and v.codeDpt2=d2.codeDpt)
    or
      (d1.codeDpt > d2.codeDpt and v.codeDpt1 = d1.codeDpt2 and v.codeDpt2=d1.codeDpt)
     )
\end{minted}

Seconde solution avec une union

\begin{minted}{sql}
select r1.intitulé as nomRégion1, r2.intitulé as nomRégion2
from  Région as r1, Région as r2, 
        Département as d1, Département as d2, Voisins as v
where d1.codeRégion = r1.codeRégion
and d2.codeRégion = r2.codeRégion
and d1.codeDpt < d2.codeDpt 
and v.codeDpt1 = d1.codeDpt1 
and v.codeDpt2 = d2.codeDpt
union
select r1.intitulé as nomRégion1, r2.intitulé as nomRégion2
from  Région as r1, Région as r2, 
        Département as d1, Département as d2, Voisins as v
where d1.codeRégion = r1.codeRégion
and d2.codeRégion = r2.codeRégion
and d1.codeDpt > d2.codeDpt 
and v.codeDpt1 = d2.codeDpt1 
and v.codeDpt2 = d1.codeDpt
\end{minted}

\end{comment}
\item 
\begin{minted}{sql}
select *
from  Région as r
where not exists (select* from Département as d
                  where r.codeRégion = d1.codeRégion)
\end{minted}

\item 
\begin{minted}{sql}
select *
from  Personne as p
where not exists (select* from Département as d
                  where d.idPrésident = p.idPrésident)
and not exists (select* from Région as d
                  where r.idPrésident = p.idPrésident)
\end{minted}

\item 
\begin{minted}{sql}
select r.intitulé, count(d.codeDpt), ssum(population) as population
from  Région as r, Département as d
where r.codeRégion = d1.codeRégion
group by r.codeRégion
\end{minted}


\item Si on veut prendre en compte les régions sans département, il faut
faire une jointure externe
\begin{minted}{sql}
select r.intitulé, count(d.codeDpt), sum(population) as population
from  Région as r left join Département as d on r.codeRégion = d1.codeRégion
group by r.codeRégion
\end{minted}
Les régions sans département apparaissent dans le résultat de la jointure avec une
valeur à \texttt{null} pour la population, et le \texttt{count} appliqué à \texttt{null}
renvoie \texttt{0} (tandis que le \texttt{sum} appliqué à \texttt{null}
renvoie \texttt{null}).

\end{enumerate}
\end{solution}


\section{Modèle relationnel ({\bf 3 pts})}

Soit les attributs TGVER avec les dépendances fonctionnelles suivantes\,:
$TG \to E$; $R \to G$; $E \to G$.

\begin{enumerate}
  \item Quelle est la clé\,?
  \item Quel est le résultat de l'algorithme de normalisation\,?
  \item Ce résultat est-il en troisième forme normale\,?
\end{enumerate}


\begin{solution}
Tous les attributs qui n'apparaissent pas à droite sont nécessairement
dans la clé, donc T, V  et R. On vérifie aisément que TVR est une clé
(la seule).

On a donc les relations (TGE) et (RG), auxquelles il faut ajouter la clé
(TVR). 

Le résultat est en troisième forme normale, si l'on admet une petite extension\,:
dans la relation TGE, la DF E -> G n'a pas la clé pour partie gauche. Mais on ne peut 
pas décomposer plus sans perdre d'information. C'est un cas (très rare en pratique) où il
faut admettre une définition de lea 3FN un peu plus compliquée que celle donnée en cours.


\end{solution}
\end{document}

